% ECONOMICS_PROBLEM: {"id": "O31-510", "title": "Entrant versus incumbent innovation", "jel_code": "O31", "source_id": "OP-0266"}
% FIRST_POSTED: 2026-10-09
% ATTEMPT_STATUS: partial
% ENTRY_KIND: research_attempt
\documentclass[11pt,a4paper]{article}
\usepackage[T1]{fontenc}
\usepackage[utf8]{inputenc}
\usepackage[margin=27mm]{geometry}
\usepackage{amsmath,amssymb,amsthm,mathtools}
\usepackage[hidelinks]{hyperref}
\setlength{\parindent}{0pt}
\setlength{\parskip}{5pt}
\newtheorem{theorem}{Theorem}[section]
\newtheorem{proposition}[theorem]{Proposition}
\newtheorem{lemma}[theorem]{Lemma}
\newcommand{\R}{\mathbb R}
\newcommand{\E}{\mathbb E}
\DeclareMathOperator{\Var}{Var}
\DeclareMathOperator{\Cov}{Cov}
\begin{document}
\section*{Entrant pressure and incumbent innovation: explicit mechanisms and an equilibrium reversal}
\section{Scope}
The derivative of incumbent research with respect to entrant displacement pressure is not universally positive or negative. This attempt supplies one strictly concave model with a tractable sign threshold, a research-labor equilibrium that can reverse a positive partial response, and sharp linear instrumental-variable sensitivity bounds. It does not estimate the causal response curve of actual firms. The catalogue already states the implicit-function identity; the purpose here is to identify primitives and information requirements that give it economic content.

Let incumbent own-product research intensity be $r\ge0$ and entrant displacement hazard $\lambda\ge0$. Fix discount rate $\rho>0$ and a research cost coefficient $w>0$. Suppose expected discounted net profit is
\[
 V(r,\lambda;w)=\frac{(a+b\lambda)r-\tfrac12cr^2}{\rho+\lambda}
                    -\tfrac12wr^2,                      \tag{1}
\]
where $a>0$, $b\ge0$, $c>0$. The denominator captures replacement risk. The numerator allows entrant pressure to increase the profit reward to a quality improvement; this is a specified escape incentive, not an inference from entry counts. The objective is strictly concave and goes to minus infinity as $r$ grows, so its unique optimum is interior and finite:
\[
 r^*(\lambda;w)=\frac{a+b\lambda}{c+w(\rho+\lambda)}.        \tag{2}
\]
The coefficient $w$ may represent the price scale of increasing marginal research inputs, rather than a literal constant wage per unit of $r$.

\section{An explicit strategic sign threshold}
Differentiating (2) gives
\[
 \frac{\partial r^*}{\partial\lambda}
 =\frac{b(c+w\rho)-wa}{[c+w(\rho+\lambda)]^2}.             \tag{3}
\]
The numerator is independent of $\lambda$ in this benchmark. Hence replacement risk dominates for $b(c+w\rho)<wa$, escape incentives dominate for the opposite strict inequality, and the response is flat at equality. Strict concavity verifies that this derivative describes the global optimum, not just a stationary point.

For $a=c=w=\rho=1$, setting $b=0$ gives $r^*=1/(2+\lambda)$ and substitution. Setting $b=1$ gives $r^*=(1+\lambda)/(2+\lambda)$ and complementarity, with derivative $1/(2+\lambda)^2$. These examples share the same replacement denominator and research-cost curvature. A falling research response is therefore not implied by displacement risk alone once the quality payoff changes with entrant pressure. Conversely, counting more entrants does not establish that the positive $b$ channel exists.

A finite-difference version does not require estimating an infinitesimal derivative. For $\lambda_1>\lambda_0$, direct subtraction in (2) yields
\[
 r^*(\lambda_1;w)-r^*(\lambda_0;w)
 =\frac{[b(c+w\rho)-wa](\lambda_1-\lambda_0)}
 {[c+w(\rho+\lambda_1)][c+w(\rho+\lambda_0)]}.             \tag{4}
\]
If feasible research is capped at $\bar r$, the best response is $\min\{\bar r,r^*\}$. Its monotonicity follows from (3), while strict derivatives can become zero at the cap and fail to exist at its boundary. A smooth interior formula should not be fitted across such a kink without acknowledging the constraint.

\section{A complete research-input price feedback}
Now endogenize the research cost scale through a reduced-form research-input market:
\[
 w=w_0+\chi r+\psi\lambda,\qquad
 w_0>0,\quad\chi,\psi\ge0.                              \tag{5}
\]
An individual incumbent treats $w$ as given; (5) represents the aggregate inverse supply evaluated at symmetric incumbent demand $r$ plus entrant demand proportional to $\lambda$. Let $D=\rho+\lambda$. Combining (2) and (5), the equilibrium research intensity is the nonnegative root of
\[
 F(r,\lambda):=\chi D r^2+[c+w_0D+\psi\lambda D]r
                         -(a+b\lambda)=0.                \tag{6}
\]
For every $\lambda\ge0$, $F(0,\lambda)<0$, $F(r,\lambda)\to\infty$, and
$F_r=2\chi Dr+c+w_0D+\psi\lambda D>0$. Thus the equilibrium exists and is unique. Implicit differentiation gives
\[
 \frac{dr}{d\lambda}=
 \frac{b-\chi r^2-[w_0+\psi(\rho+2\lambda)]r}
 {2\chi Dr+c+w_0D+\psi\lambda D}.                         \tag{7}
\]
The denominator is positive. The new numerator includes input-market congestion and entrant demand for those inputs, neither of which is in the fixed-$w$ derivative (3).

Take $a=b=c=\rho=w_0=1$, $\chi=0$, $\psi=3$. At $\lambda=0$ the equilibrium is $r=1/2$, $w=1$. Holding $w$ fixed gives the positive partial derivative $1/4$ from (3). Along equilibrium, (7) gives
$[1-(1+3)(1/2)]/2=-1/2$. Thus the same point has strategic complementarity at fixed research prices and a negative total response when entrant pressure raises input costs. The inverse supply equation is a declared market primitive, so this is an internally closed benchmark rather than a partial derivative mislabeled as general equilibrium. Endogenous product demand and finance would add further equations and need their own existence argument.

\section{Instrumental variation with a bounded exclusion violation}
Consider a separate local linear empirical specification in a prespecified technological-distance group:
\[
 R=\alpha+\beta\Lambda+U,
\]
where $R$ is incumbent own-product research and $\Lambda$ is own-line displacement hazard. Center a candidate instrument $Z$ and suppose the observable population moments are
$q=E[Z\Lambda]\ne0$ and $s=E[ZR]$. An exact exclusion sets $E[ZU]=0$ and gives $\beta=s/q$. If instead $|E[ZU]|\le\Delta$, then the exact moment-compatible interval is
\[
 \beta\in\left[\frac sq-\frac{\Delta}{|q|},\quad
                    \frac sq+\frac{\Delta}{|q|}\right].   \tag{8}
\]
Necessity follows from $s-\beta q=E[ZU]$. For any retained $\beta$, define the residual from the displayed structural equation using the observed $(R,\Lambda,Z)$ and choose the intercept to match the mean. Its moment then obeys the restriction, proving sharpness within this linear moment class. The construction does not prove that $\beta$ is a causal derivative outside the maintained linear model.

For example $q=0.1,s=0.02,\Delta=0.03$ gives $[-0.1,0.5]$; a positive naive ratio $0.2$ does not establish complementarity under that exclusion allowance. An instrument that also changes incumbent demand or research productivity can violate the intended fixed-price strategic interpretation even if it predicts entry strongly.

Sampling inference can avoid division by a weak estimated first stage: retain every candidate $\beta$ for which a simultaneous confidence region for $(s,q)$ contains some pair satisfying $|s-\beta q|\le\Delta$. On the region's coverage event this projection covers the true compatible slopes. If the region includes $q=0$ and $|s|\le\Delta$, the retained slope set can be unbounded. Imposing a finite interval solely for plotting would conceal weak identification. Tail conditions or clustered sampling assumptions are needed to construct the moment region; none is supplied by the algebra itself.

\section{Measurement and unresolved empirical content}
An incumbent's acquisition expenditure is not its own-product research intensity, and total new-firm entry is not necessarily a hazard of displacement in its product line. A credible application would match research inputs, innovations in existing products, rival quality advances, and input prices within technological-distance groups. It would identify the fixed-price response separately from wage and financing feedback, and test the maintained exclusion with relevant institutional evidence rather than with a first-stage coefficient alone.

The benchmark resolves conditional signs and provides a concrete reversal, but the actual sign function across technologies remains an empirical question. The growth-accounting source measures contributions of different innovation types; accounting shares do not identify the strategic causal derivative in (3) or (7). No universal complementarity or substitutability claim is made.

\section{Completion Estimate}
50\%

This is a post hoc, heuristic estimate for the full catalogue target.
The attempt derives replacement-escape sign thresholds, an equilibrium research-input-price reversal, and sharp local linear instrument-exclusion sensitivity bounds. The full causal strategic-response function still requires matched incumbent own-product innovation and entrant displacement data, credible fixed-price exclusions, and validated heterogeneity across technological distance.

\begin{thebibliography}{9}
\bibitem{kl} P. J. Klenow and H. Li (2021), \emph{Innovative Growth Accounting}, NBER Macroeconomics Annual 35. \url{https://www.journals.uchicago.edu/doi/10.1086/712325}. Working-paper record: \url{https://www.nber.org/papers/w27015}. The source motivates the relationship between incumbent innovation and creative destruction; the profit specification and equilibrium reversal above are explicit illustrative primitives.
\end{thebibliography}

\end{document}
